\documentclass{article}
\usepackage{amsmath,amssymb}
\usepackage{hyperref}
\usepackage{fullpage}

\title{Issue \#1622: the degree-zero branch of low-degree pasting}
\author{MIPStarRE formalization notes}
\date{2026-05-14}

\begin{document}
\maketitle

\section{Source statement}

The source theorem is \texttt{references/ldt-paper/ld-pasting.tex},
lines 12--50, labelled \texttt{thm:ld-pasting}.  It assumes an
$(\varepsilon,\delta,\gamma)$-good symmetric strategy for the
$(m+1,q,d)$ low individual degree test, a family of projective
submeasurements
\[
  \{G^x\}_{x\in\mathbb F_q}\subseteq \mathrm{PolySub}(m,q,d),
\]
the four displayed completeness, consistency, strong self-consistency, and
boundedness hypotheses, and an integer \(k\geq 400md\).  It concludes that
there is a pasted measurement \(H\in\mathrm{PolyMeas}(m+1,q,d)\) satisfying the
final point-consistency estimate with error
\[
  \sigma=\kappa\left(1+\frac{1}{100m}\right)+2\nu+
  \exp\left(-\frac{k}{80000m^2}\right),
\]
where
\[
  \nu=100k^2m\left(\varepsilon^{1/32}+\delta^{1/32}
    +\gamma^{1/32}+\zeta^{1/32}+(d/q)^{1/32}\right).
\]

The theorem does not state \(0<d\).  The following paragraph of the paper,
lines 52--55, only says that the proof may assume that
\(\varepsilon,\delta,\gamma,\zeta,d/q\leq 1\), since the bound is trivial when
one of these quantities is at least \(1\).

\section{Lean status}

The source-facing Lean theorem is
\[
  \texttt{MIPStarRE.LDT.Pasting.ldPasting}.
\]
Its public statement keeps the paper's unrestricted boundary.  The restricted
nontrivial construction theorem
\[
  \texttt{MIPStarRE.LDT.Pasting.ldPastingNontrivial}
\]
assumes the proof-regime hypotheses
\[
  \gamma\leq 1,\qquad \zeta\leq 1,\qquad d\leq q,\qquad 0<d,\qquad 1\leq k.
\]
The unrestricted theorem dispatches the large-\(\gamma\), large-\(\zeta\),
large-\(d/q\), \(k=0\), and nontrivial positive-degree branches.  The remaining
branch is
\[
  d=0,\qquad 1\leq k,
\]
recorded as
\[
  \texttt{MIPStarRE.LDT.Pasting.ldPastingDegreeZeroBranch}.
\]
That declaration currently contains the direct tracked \texttt{sorry} for
issue \#1622.

\section{Why the existing proof does not specialize}

The existing nontrivial proof uses the sampled distinct-tuple construction from
\texttt{ld-pasting.tex}.  In the proof of \texttt{lem:h-b-consistency}, lines
1073--1087, and again in the proof of \texttt{lem:over-all-outcomes}, lines
1192--1199 and 1276--1286, the argument passes between independent sampled
heights and distinct sampled heights.  Each such passage costs a
distinctness error of order \(k^2/q\).

For \(0<d\), the formal proof absorbs this cost through the displayed
\((d/q)^{1/32}\) term.  Concretely, the scalar step uses \(1/q\leq d/q\) and
then \(d/q\leq (d/q)^{1/32}\) in the small-error regime.  This is exactly the
role of the Lean hypothesis \(0<d\) in the scalar aggregation lemmas such as
\[
  \texttt{one\_div\_q\_le\_rpow\_degreeRatio}
\]
and its downstream uses in the pasting proof.

When \(d=0\), this absorption is unavailable: the term
\((d/q)^{1/32}\) is zero.  Thus the positive-degree proof cannot be reused by
specialization, and adding \(0<d\) to \texttt{ldPasting} would change the source
theorem rather than prove the missing branch.

\section{Obligation}

The remaining construction theorem should prove the degree-zero branch from the
paper hypotheses:
\[
  d=0,\qquad 1\leq k,\qquad k\geq 400md.
\]
Mathematically, one expects a proof which uses the fact that a degree-zero
polynomial is constant and obtains point-consistency through the line-test and
slice-consistency hypotheses, without paying a distinct-tuple loss that must be
absorbed by \(d/q\).  The precise formal target is
\[
  \texttt{MIPStarRE.LDT.Pasting.ldPastingDegreeZeroBranch}.
\]

Until that theorem is proved, the project should keep the source-facing theorem
\texttt{ldPasting} visible and should keep the missing work as the tracked
\texttt{sorry} in this branch.  It should not add \(0<d\), a bridge hypothesis,
a residual package, or any other non-paper assumption to
\texttt{thm:ld-pasting}.

\end{document}
